% Section F: practice (LP21-LP22).
\section{In practice}

\begin{frame}{Lever arm: where you mount the IMU matters \hfill \lp{22}}
  \begin{columns}[T]
    \begin{column}{0.48\textwidth}
      \small
      An IMU at offset $\vect{r}$ from the rotation centre measures
      \[
        \vect{f}_{\mathrm{IMU}} = \vect{f}_{\mathrm{c}} + \boldsymbol{\omega}\times(\boldsymbol{\omega}\times\vect{r}) + \dot{\boldsymbol{\omega}}\times\vect{r}
      \]
      \begin{itemize}
        \item Centripetal $\omega^2 r$ points toward the axis: a spinning robot ``leans''.
        \item Tangential $\dot\omega\, r$: every yaw acceleration kicks the reading.
        \item Fix: mount near the centre, or compensate using the gyro rate and a calibrated $\vect{r}$. VIO/LIO pipelines estimate $\vect{r}$ online.
      \end{itemize}
    \end{column}
    \begin{column}{0.5\textwidth}
      \includegraphics[width=\linewidth]{lever_arm_tilt.pdf}
    \end{column}
  \end{columns}
  \tryit{accelerometer}{lever\_arm\_spin}
\end{frame}
